% GENERATED from content/data/enrichment.json by mxbuild.enrichment; do not edit.
\mxunit{enrichment-inequalities-notes}{Inégalités et cas d’égalité}
\mxsubtitle{Lycée · optimisation et olympiades \textperiodcentered{} Leçon et méthode \textperiodcentered{} 11}

\begin{prerequisites}
Carrés, factorisation et racines carrées.
\end{prerequisites}

\begin{keypoint}[Idées et méthode]
Pour tout réel $u$, $u^{2}\geq 0$, avec égalité exactement lorsque $u=0$. Beaucoup d’inégalités se ramènent à cette observation.

Pour $a,b\geq 0$, $(\sqrt{a}-\sqrt{b})^{2}\geq 0$ donne $a+b\geq 2\sqrt{ab}$. C’est l’inégalité entre les moyennes arithmétique et géométrique. L’égalité a lieu si et seulement si $a=b$.

Prouver une borne ne suffit pas à annoncer un maximum : il faut vérifier qu’un cas d’égalité autorisé l’atteint. Les contraintes de positivité et de domaine font partie du raisonnement.
\end{keypoint}

\begin{example}
Si $a+b=10$ avec $a,b\geq 0$, alors $(a-b)^{2}\geq 0$ donne $100-4ab\geq 0$, donc $ab\leq 25$. Le couple $a=b=5$ atteint 25 : c’est bien le maximum.
\end{example}

\begin{activity}[À explorer]
Tracer $f(x)=x(10-x)$ pour $0\leq x\leq 10$ et déplacer le point. Retrouver ensuite le maximum sans utiliser le dessin.
\end{activity}

